\documentclass[12pt]{article}

\usepackage[margin=1.12in]{geometry}
\usepackage{amsmath,amssymb,amsthm,mathtools}
\usepackage{booktabs}
\usepackage{hyperref}

\newtheorem{theorem}{Theorem}
\newtheorem{proposition}{Proposition}
\theoremstyle{remark}
\newtheorem{remark}{Remark}

\newcommand{\A}{\mathcal A}
\newcommand{\B}{\mathcal B}
\newcommand{\C}{\mathbb C}
\newcommand{\Hcode}{\mathcal H_{\mathrm{code}}}
\newcommand{\M}{\mathcal M}
\newcommand{\Tr}{\operatorname{Tr}}
\newcommand{\id}{\operatorname{id}}

\setlength{\emergencystretch}{3em}
\hypersetup{
  pdftitle={Algebra-Relative OAQEC/JLMS Reconstruction Does Not Certify Maximality},
  pdfsubject={Operator-algebra quantum error correction, JLMS, finite centers, and maximality},
  pdfkeywords={OAQEC, JLMS, entanglement wedge reconstruction, finite centers, algebra-relative recovery}
}

\title{Algebra-Relative OAQEC/JLMS Reconstruction\\
Does Not Certify Maximality}
\author{}
\date{Live review note of June 20, 2026}

\begin{document}
\maketitle

\begin{abstract}
Let \(\M\subseteq\B(\Hcode)\) be a finite-dimensional von Neumann algebra
reconstructable on a boundary region or boundary algebra in the usual
operator-algebra quantum-error-correction sense.  For every \(q\geq 2\), the
finite central direct-sum extension
\[
  \Hcode^{(q)}=\bigoplus_{j=1}^{q}\Hcode,\qquad
  \M_q=\M\otimes \C^q\simeq \bigoplus_{j=1}^{q}\M
\]
is reconstructable blockwise whenever the boundary representation and recovery
are extended blockwise.  A quotient-visible record obtained by composing with
a central character \(\pi_k:\C^q\to\C\) sees only the selected summand:
\[
  Q_k=\id_{\M}\otimes \pi_k:\M_q\to \M.
\]
Thus OAQEC/JLMS data certify reconstruction of the supplied algebra.  They do
not, by themselves, certify that the supplied algebra is the full logical
algebra against hidden finite central refinements.  The obstruction disappears
only after adding a maximality premise, complementary recovery/full-algebra
identification, center-separating states or operations, or a physical rule
excluding the finite center.
\end{abstract}

\paragraph{Status.}
This is non-frozen external-review metadata for Team 1.  It is intended as a
small claim that a holographic-QEC referee can classify as known, false,
missing-hypothesis, or useful.  It does not modify any frozen Team 1
manuscript, Lean archive, public mirror manifest, or SHA-256 package.

\section{Setup}

Let \(\Hcode\) be finite-dimensional and let
\(\M\subseteq\B(\Hcode)\) be a finite-dimensional von Neumann algebra.  The
finite-dimensional structure theorem writes
\[
  \M \simeq \bigoplus_{\alpha}
  \bigl(\B(\mathcal H_{a_\alpha})\otimes I_{\bar a_\alpha}\bigr).
\]
In operator-algebra quantum error correction, a boundary region or boundary
algebra reconstructs \(\M\) when every \(x\in\M\) has a boundary
representative acting identically on the code, equivalently when the relevant
erasure/recovery, commutant, and relative-entropy tests hold for the declared
algebra.

Fix \(q\geq 2\).  Define
\[
  \Hcode^{(q)}=\bigoplus_{j=1}^{q}\Hcode,\qquad
  \M_q=\M\otimes \C^q\simeq\bigoplus_{j=1}^{q}\M.
\]
Let \(p_j\) be the minimal central idempotents of \(\C^q\).  For a selected
character \(\pi_k:\C^q\to\C\), set
\[
  Q_k=\id_{\M}\otimes \pi_k:\M_q\to\M.
\]

\begin{theorem}[finite central direct-sum extension]
\label{thm:oaqec-direct-sum}
Suppose \(\M\) is reconstructable on a boundary region or algebra \(A\) in the
finite-dimensional OAQEC sense.  If the code, boundary representation, and
recovery map are extended blockwise to \(\Hcode^{(q)}\), then
\(\M_q=\M\otimes\C^q\) is reconstructable blockwise on \(A\).  However, every
declared reconstruction, recovery, modular, entropy, or JLMS record that
factors through \(Q_k\) is insensitive to the nonselected central supports
\(\sum_{j\ne k}p_j\) and to the hidden central cardinality \(q\).
\end{theorem}

\begin{proof}
Let \(R_A:\M\to\B(\mathcal H_A)\) denote a boundary representation of the
declared algebra.  Define
\[
  R_A^{(q)}(x_1,\ldots,x_q)
  =
  \bigoplus_{j=1}^{q} R_A(x_j)
\]
on the blockwise boundary representation space.  Since each component
\(R_A(x_j)\) acts on the corresponding copy of the code as \(x_j\), the direct
sum acts on \(\Hcode^{(q)}\) as \((x_1,\ldots,x_q)\).  The same componentwise
construction extends any recovery channel or algebra homomorphism used to
state the OAQEC condition.  Therefore \(\M_q\) is reconstructable blockwise.

Now suppose a record factors through \(Q_k\).  If
\[
  (x_1,\ldots,x_q),(y_1,\ldots,y_q)\in\M_q
  \quad\text{with}\quad x_k=y_k,
\]
then \(Q_k(x_1,\ldots,x_q)=Q_k(y_1,\ldots,y_q)\).  Consequently every
boundary representative, recovery task, modular comparison, or relative
entropy equality evaluated after the quotient has the same value for the two
tuples, even if their nonselected components or central supports differ.  The
quotient-visible data therefore cannot determine the omitted finite central
fiber or the number \(q\).
\end{proof}

\section{Relation To JLMS}

The JLMS statement compares boundary and bulk relative entropies for the bulk
algebra or entanglement-wedge algebra under discussion.  In the
operator-algebra QEC formulation, the theorem is algebra-relative: the
recovered object is the stated von Neumann algebra, possibly with center.  The
direct-sum extension above leaves that algebra-relative statement intact.

\begin{proposition}[quotient-visible JLMS nonidentifiability]
\label{prop:jlms-quotient}
Consider any family of code states, modular-flow tests, or relative-entropy
comparisons whose declared bulk algebra is \(Q_k(\M_q)\simeq\M\), or whose
observed records factor through \(Q_k\).  The resulting JLMS/OAQEC data do not
distinguish \(\M\) from the selected quotient of \(\M\otimes\C^q\), and hence
do not certify that no hidden finite central direct-sum refinement was present.
\end{proposition}

\begin{proof}
Relative entropy and modular data in the declared quotient are functions of
the selected block.  Changes to nonselected blocks change neither the
quotient state nor the quotient algebra element.  Thus the same
quotient-visible JLMS/OAQEC record is compatible with multiple values of
\(q\) and with multiple hidden central supports.
\end{proof}

\begin{remark}
If instead the state family includes all central sectors and their classical
probabilities, full direct-sum relative entropy contains the usual classical
relative-entropy term between central distributions.  That is not a
counterexample to the obstruction; it is a center-separating positive exit.
\end{remark}

\section{Positive Exits}

The obstruction is intentionally conditional.  It is killed, or demoted to a
routine application, by any hypothesis that supplies finite-center
accountability:

\begin{center}
\begin{tabular}{p{0.33\linewidth}p{0.53\linewidth}}
\toprule
Exit premise & Effect \\
\midrule
Full logical algebra stipulated & If \(\M\) is defined to be the complete
logical algebra of the code, then \(\M\otimes\C^q\) is a different code
algebra, not a hidden refinement of the same target. \\
\addlinespace
Complementary recovery/maximality & If the theorem being invoked proves a
unique maximal algebra from both region and complement data, then omitted
central summands are excluded by hypothesis or theorem. \\
\addlinespace
Center-separating states or operations & If the comparison includes states,
operations, or boundary readouts that distinguish the central projections
\(p_j\), then the data no longer factor through \(Q_k\). \\
\addlinespace
Physical exclusion & If the holographic model forbids such finite classical
superselection labels, the extra premise should be named. \\
\bottomrule
\end{tabular}
\end{center}

\section{Referee Question}

The classification question is precise:

\begin{quote}
Does the standard OAQEC/JLMS hypothesis being used already imply that the
supplied algebra is maximal as the full logical algebra, or that all finite
central sectors are measured, included, or physically excluded?  If so, name
the theorem or axiom.  If not, algebra-relative reconstruction does not by
itself certify finite-center accountability.
\end{quote}

The desired external response is not encouragement.  It is one of: known
finite-dimensional OAQEC fact, false statement, missing positive hypothesis,
physically irrelevant comparison class, or useful expository no-go.

\begin{thebibliography}{99}

\bibitem{JLMS}
D. L. Jafferis, A. Lewkowycz, J. Maldacena, and S. J. Suh,
Relative entropy equals bulk relative entropy,
\emph{Journal of High Energy Physics} \textbf{2016}, 004 (2016);
\texttt{arXiv:1512.06431}.

\bibitem{DHW}
X. Dong, D. Harlow, and A. C. Wall,
Reconstruction of bulk operators within the entanglement wedge in gauge-gravity
duality,
\emph{Physical Review Letters} \textbf{117}, 021601 (2016);
\texttt{arXiv:1601.05416}.

\bibitem{Harlow}
D. Harlow,
The Ryu-Takayanagi formula from quantum error correction,
\emph{Communications in Mathematical Physics} \textbf{354} (2017), 865--912;
\texttt{arXiv:1607.03901}.

\bibitem{KamalPenington}
A. A. Kamal and G. Penington,
The Ryu-Takayanagi formula from quantum error correction: an algebraic
treatment of the boundary CFT,
\texttt{arXiv:1912.02240}.

\end{thebibliography}

\end{document}
